\documentclass{homework}
\course{Math 4182H}
\author{Jim Fowler}
\hwtitle{Honors Analysis II}
\input{preamble}

\newcommand{\area}{\operatorname{area}}
\newcommand{\vol}{\operatorname{vol}}
\newcommand{\Jac}{\operatorname{Jac}}

\renewcommand*{\thefootnote}{\fnsymbol{footnote}}

\begin{document}
\maketitle

\begin{inspiration}
Le leggi della logica non bastano a creare l'opera del matematico.\footnote{``The laws of logic are not enough to create the work of the mathematician.''}
\byline{Guido Fubini}
\end{inspiration}

\section{Terminology}

\begin{problem}
  Let $S\subset \R^2$ be a bounded region.

  What does it mean to say $S$ is a \textbf{type I region}?
  A \textbf{type II region}?
\end{problem}

\section{Numericals}

\begin{problem}\label{calabi}%
Show that
\[
\lim_{\epsilon \to 0^{+}} \int_{0}^{1-\epsilon}\int_{0}^{1-\epsilon}\frac{dx\,dy}{1-x^{2}y^{2}}
=\sum_{n=0}^{\infty}\frac{1}{(2n+1)^{2}}
=\frac{3}{4}\sum_{n=1}^{\infty}\frac{1}{n^{2}}.
\]
\end{problem}

\begin{problem}
Compute the integral and sum appearing in \ref{calabi}.

\textit{Hint:} Let $T$ be the right triangle with vertices $(0,0)$ and $(0,\pi/2)$ and $(\pi/2,0)$. Use the transformation $f\colon T\to\mathbb{R}^{2}$ defined by
\[
f(u,v)=\Bigl(\frac{\sin u}{\cos v},\,\frac{\sin v}{\cos u}\Bigr)
\]
to verify
\[
\int_{0}^{1}\int_{0}^{1}\frac{1}{1-x^{2}y^{2}} \, dx \, dy
=\iint_{T}1\,du\,dv
\]
and then to ultimately compute $\displaystyle\sum_{n=1}^{\infty}\frac{1}{n^{2}}$.
\end{problem}

\section{Exploration}

\begin{problem}\label{product-formula}%
  Suppose $f : [a,b] \to \R$ and $g : [c,d] \to \R$ are continuous. Consider the rectangle $S = [a,b] \times [c,d]$ and define $h : S \to \R$ by $h(x,y) = f(x) \, g(y)$, and then relate $\iint_S h \, dA$ to $\int_a^b f(x) \, dx$ and $\int_c^d g(y) \, dy$.
\end{problem}

\begin{problem}
Define
\[
I := \int_{-\infty}^{\infty} e^{-x^2}\,dx.
\]
Use \ref{product-formula} to show
\[I^2
=
\int_{\mathbb{R}^2} e^{-(x^2+y^2)}\,dx\,dy,
\]
and then switch to polar coordinates to prove $I = \sqrt{\pi}$.
\end{problem}

\begin{problem}
Define the $\Gamma$-function for $s>0$ by
\[
\Gamma(s):=\int_{0}^{\infty} t^{s-1}e^{-t}\,dt.
\]
For $n\in\mathbb{N}$, set
\[
G_n := \int_{\mathbb{R}^n} e^{-\|x\|^2}\,dx.
\]
Define $f(r) = e^{-r^2}$, and 
explain why there is a constant $A_{n-1}$ (the surface area of the unit
$(n-1)$-sphere) such that 
  \[
  \int_{\mathbb{R}^n} f(\|x\|)\,dx \;=\; \int_{0}^{\infty} f(r)\,A_{n-1}\,r^{n-1}\,dr.
  \]
Therefore
  \[
  G_n = \pi^{n/2} = A_{n-1}\int_{0}^{\infty} r^{n-1}e^{-r^2}\,dr.
  \]
Compute the integral
  \[
  \int_{0}^{\infty} r^{n-1}e^{-r^2}\,dr
  \]
  in terms of the $\Gamma$-function, and deduce a formula for
  $A_{n-1}$.

Finally, compute the volume $V_n$ of the unit $n$-ball
  \[
  B^n := \{x\in\mathbb{R}^n : \|x\|\le 1\}.
  \]
\end{problem}



\begin{problem}
Let $R\subset \mathbb{R}^2$ be a rectangle with axis-parallel sides. Suppose $R$ is partitioned into finitely many axis-parallel subrectangles
\[
R = R_1 \sqcup R_2 \sqcup \cdots \sqcup R_n
\]
whose interiors are pairwise disjoint and whose union is $R$. Assume that for each $i$, the rectangle $R_i$ has \emph{at least one} side length that is an integer.

Prove that $R$ must also have at least one side length that is an integer.

\textit{Hint:} Consider the function $f(x,y)=e^{2\pi i(x+y)}$ and the integral
\[
I(S) \;:=\; \int\!\!\int_{S} e^{2\pi i(x+y)}\,dx\,dy
\]
for an axis-parallel rectangle $S=[a,b]\times[c,d]$ and 
show that $I(S)=0$ if and only if at least one of the side lengths $b-a$ or $d-c$ is an integer.
Then use additivity of the integral over a partition, i.e., \(I(R)=\sum_{i=1}^n I(R_i)\).
\end{problem}

\section{Prove or Disprove and Salvage if Possible}

\begin{problem}\label{fubini-no-hypotheses}%
  If $f$ is Riemann integrable on $R=[0,1]\times[0,1]$, then for every $y\in[0,1]$ the function
  $x\mapsto f(x,y)$ is Riemann integrable on $[0,1]$.
\end{problem}

\begin{problem}\label{swap-always}%
  If the two iterated (Riemann!) integrals
  \[
    \int_0^1\left(\int_0^1 f(x,y)\,dx\right)\,dy
    \qquad\text{and}\qquad
    \int_0^1\left(\int_0^1 f(x,y)\,dy\right)\,dx
  \]
  both exist, then they are equal.
\end{problem}

\end{document}
